\section{The augmented representation, and the sign convention}
\label{app:representation}

\subsection{Assumptions under which \eqref{eq:augmented} reduces exactly}

Write the augmented weights as $W_r=[\hat{w}_r; a_r]$ with a Gaussian belief of
variance $v_a$ on the extra coordinate. The general compatibility score is
%
\begin{equation}
  \Hf^{\text{aug}}
  = \frac{\hw + \sigma_e\,a_r z_r s/\gamma_C}{\sqrt{1+z_r v_a/\gamma_C^{2}}}
  = \hw + \frac{D}{\gamma_C} - \frac{\hw z_r v_a}{2\gamma_C^{2}} + O(v_a^{2}),
\end{equation}
%
so \eqref{eq:disc-field} is recovered exactly under four assumptions, which we
state rather than leave implicit.

\textbf{A1 (deterministic coordinate, $v_a=0$).} The class coordinate is part of
the representation, not a learned parameter, so it carries no posterior variance.
This is what leaves $\gamma_C$ untouched and suppresses any covariance update
along that direction. It is the only assumption needed for exactness; relaxing it
gives the $O(v_a/\gamma_C^2)$ correction above.

\textbf{A2 (fixed prejudice).} The weight on the coordinate is the hyperparameter
$d$ and is not updated. Relaxing A2 gives a \emph{learned-prejudice} variant in
which $a_r \leftarrow a_r + (\Fw/\gamma)v_a z_r s$, so the strength of the bias
would itself adapt to how well it predicted. That is a natural extension and we
do not study it.

\textbf{A3 (interaction-level, not issue-level).} The coordinate attaches to the
message, because $\kappa_e$ is a property of the source and not of the issue.
This is what keeps the shift from being multiplied by $\sigma_e$ and divided by
$\gamma_C$ --- either of which would make the bias flip with the emitter's stated
opinion, or decay as the receiver grew confident.

\textbf{A4 (bounded binary feature).} $s\in\{0,\pm1\}$, so the shift is bounded
by $|d|$ and there are exactly six inequivalent fillings of the $2\times2$
matrix, tabulated below.

\subsection{The six discrimination patterns}

Rows index the receiver's class, columns the emitter's; $+d$ is the tolerant
entry and $-d$ the intolerant one.

\newcommand{\dmat}[4]{\left(\begin{smallmatrix}#1&#2\\#3&#4\end{smallmatrix}\right)}
\begin{table}[h]
\centering\small
\begin{tabular}{cccc}
\toprule
case & $D/d$ & who discriminates & how \\
\midrule
1 & $\dmat{+1}{0}{0}{0}$    & $A$ only & in-group favouritism \\
2 & $\dmat{0}{-1}{0}{0}$    & $A$ only & out-group hostility \\
3 & $\dmat{+1}{-1}{0}{0}$   & $A$ only & both \\
4 & $\dmat{+1}{0}{0}{+1}$   & both     & in-group favouritism \\
5 & $\dmat{0}{-1}{-1}{0}$   & both     & out-group hostility \\
6 & $\dmat{+1}{-1}{-1}{+1}$ & both     & both \\
\bottomrule
\end{tabular}
\end{table}

Case 6, used throughout the main text, is the one in which the two classes play
symmetric roles, so any asymmetry appearing in the results is spontaneous rather
than imposed. Compactly, $\Dfield=z_r d\,\kappa_r\kappa_e$.

\subsection{The sign convention, and an inconsistency in earlier presentations}
\label{app:sign}

Three statements are in play in earlier accounts of this model: that a
discriminating receiver evaluates its update at $\hw+\Dfield$; that for $d>0$ the
in-group entries of the matrix are $-d$; and that the discriminatory regime lies
at $d>0$. The first two together imply the opposite of the third.

The argument is one line of \eqref{eq:modulation}. $\Fmu$ carries the factor
$1-2\Phi(\Hf)$, which is negative whenever $\Hf>0$, and $\mu$ is \emph{distrust},
updated as $\mu\leftarrow\mu+\Fmu V/\gamma_V$. So perceived agreement drives $\mu$
down and builds trust, and perceived disagreement drives it up. A receiver adding
$+d$ therefore reads its counterpart as more agreeable than it is and grows
\emph{more trusting}: $+d$ is the tolerant entry. Assigning $-d$ to the in-group
means a society at $d>0$ distrusts its own class, which places the discriminatory
regime at $d<0$ and mirrors every map in $d$.

Any one of three repairs resolves it, and they are the same change: flip the
table so in-group entries are $+d$ (what we do, since it leaves the field
equation untouched and preserves the meaning of every published axis); or keep
the table and write the field equation with a minus; or keep both and relabel
$d<0$ as the discriminatory direction. The choice has no effect on the physics
--- the same regimes exist in the same places --- only on which half of the axis
carries which name. The accompanying code implements both readings behind a flag
and asserts that they are exact mirror images.
